%%%%%%%%%%%%%%%%%%%%%%%%%%%%%%%%%%%%%%%%%%%%%%%%%%%%%%%%%%%%%%%%%%%%%%%%%
%
%	LaTeX File for Stanford University PhD Thesis
%
%%%%%%%%%%%%%%%%%%%%%%%%%%%%%%%%%%%%%%%%%%%%%%%%%%%%%%%%%%%%%%%%%%%%%%%%%
%	Copyright 2001  by  Jung-Suk Goo    (goojs@gloworm.stanford.edu)
%%%%%%%%%%%%%%%%%%%%%%%%%%%%%%%%%%%%%%%%%%%%%%%%%%%%%%%%%%%%%%%%%%%%%%%%%
% \pagenumbering{arabic}
% \setcounter{page}{1}

\chapter{Introduction}
\label{chap:intro}

\fancypar{T}{{\sc he}} rapid evolution of wireless communications, driven by technologies such as the upcoming 6G standard\cite{fontaine2024towards}, necessitates a transformation in how Artificial Intelligence (AI) addresses challenges unattainable by traditional optimization techniques. The remarkable success of foundation models in computer vision (CLIP\cite{radford2021learning}, MAE\cite{he2022masked}) and natural language processing (GPT\cite{10500411}, BERT\cite{devlin2019bert}) has demonstrated the transformative power of learning universal representations that generalize across tasks with minimal fine-tuning. These models leverage massive unlabeled datasets through self-supervised learning to discover rich, transferable features that encode fundamental properties of their respective domains. Yet, the current AI landscape in wireless networks remains dominated by models developed and trained for isolated, task-specific applications\cite{fontaine2024towards}, leading to significant limitations in knowledge sharing, reuse, and adaptability. This dissertation addresses a fundamental question: \textit{Can we develop practical pathways toward Wireless Physical-Layer Foundation Models (WPFMs) by strategically approximating and integrating world knowledge through modular, cost-effective architectures?}

The limitations of current wireless AI systems are particularly evident in Automatic Modulation Classification (AMC)—a critical component of spectrum management, signal identification, and cognitive radio systems\cite{gao2023moe}. Conventional deep learning approaches face three major challenges: (1) \textbf{Imbalance across conditions:} Models often fail to distinguish between low and high Signal-to-Noise Ratio (SNR) conditions\cite{gao2023moe}, resulting in uneven classification performance where low SNR signals obscured by noise require specialized handling compared to high SNR signals; (2) \textbf{Data scarcity and training overhead:} Acquiring high-quality labeled radio frequency (RF) signal data demands specialized expertise and equipment, while traditional deep learning necessitates large labeled datasets\cite{kong2023transformer}, creating barriers to scalability; (3) \textbf{Lack of interpretability and generalization:} Task-specific models struggle to adapt quickly to new tasks, generalize across diverse channel characteristics, and lack capacity for human-understandable interaction\cite{gao2023moe}.

Building wireless foundation models from scratch to address these challenges is prohibitively expensive~\cite{radford2021learning}. Unlike images or text, wireless signals exhibit unique characteristics: they are inherently noisy, their structure depends on both transmitter design and channel physics, and their semantic meaning is encoded in subtle statistical patterns. These properties demand an economical approach that leverages existing world knowledge rather than massive isolated training. This dissertation presents such a solution: a modular and cost-effective framework that achieves foundation model capabilities through the strategic integration of world knowledge. Our approach validates three essential components that work synergistically: (1) a \textbf{rich feature encoder} (DenoMAE) that uses self-supervised learning to approximate world knowledge of signal representations; (2) an \textbf{intelligent reasoning engine} (LLM integration) that performs open-set classification by converting signal features into natural language prompts; and (3) an \textbf{adaptation module} (Meta-Task) that enables rapid knowledge transfer to new tasks with minimal data. By demonstrating proof-of-concept for these core components, we provide both a practical framework and a conceptual blueprint for achieving foundation model capabilities in wireless communications without prohibitive training costs.

\section{Objectives and Motivation}
\label{sec:objective}

The main objective of this dissertation is to establish foundational principles and practical methodologies for developing Wireless Physical-Layer Foundation Models (WPFMs). These models serve as universal representations of wireless signals, allowing for generalization across various modulation schemes, channel conditions, and tasks. Instead of pursuing a cost-prohibitive strategy of building massive models from scratch, this work demonstrates that the capabilities of foundation models can be achieved more economically through a strategic and modular design. This design integrates self-supervised learning, reasoning based on large language models (LLMs), and meta-learning adaptation.

The motivation behind this approach is inspired by the success of foundation models in various domains. Models such as GPT-4~\cite{achiam2023gpt} and CLIP~\cite{radford2021learning} have demonstrated that pre-training on large, diverse datasets leads to the development of powerful and transferable representations. These successes rely on three essential components: vast amounts of unlabeled data, self-supervised pretext tasks that capture domain-invariant structures, and flexible architectures that allow for adaptation to specific tasks.

In the field of wireless communications, petabytes of unlabeled signal data are generated daily. However, the unique characteristics of these signals—particularly their noisy and physics-dependent nature—pose distinct challenges. This dissertation aims to address the key challenge of developing WPFMs that can comprehend and represent these complex signals without incurring the prohibitive costs associated with isolated training, while still achieving the generalization capabilities seen in their vision and language counterparts.

To address this, the remainder of this section is structured to first articulate the long-term vision for WPFMs. It then deconstructs this vision into a set of five core research questions that guide the technical work. Finally, it outlines the specific technical approach and contributions developed in this dissertation to answer these questions.

\subsection{Vision: Wireless Foundation Models}

We envision wireless foundation models analogous to CLIP in vision or GPT in language: pre-trained on diverse, unlabeled signals to learn representations encoding fundamental wireless signal properties, then adapted to specific tasks (classification, denoising, channel estimation) with minimal fine-tuning. Such models would:

\begin{itemize}
    \item \textbf{Generalize Across Modulations:} Recognize modulation schemes never encountered during training through compositional understanding of signal properties, enabled by rich learned representations
    \item \textbf{Handle SNR Imbalance Robustly:} Employ hybrid generative-discriminative learning (DenoMAE) combined with robust discriminative classifiers (improved plug-and-play), ensuring consistent performance from -10 dB to 30+ dB
    \item \textbf{Operate Cost-Effectively:} Enable training-free or few-shot deployment through LLM integration and self-supervised pre-training, avoiding exhaustive retraining for each scenario
    \item \textbf{Adapt Rapidly with Minimal Data:} Leverage meta-learning to quickly specialize for new deployment scenarios with few labeled examples, dramatically reducing data requirements
    \item \textbf{Provide Interpretable Reasoning:} Convert signal features to natural language prompts for LLM-based classification, enabling human-understandable explanations and prompt-based optimization
    \item \textbf{Leverage World Knowledge:} Integrate pre-trained LLMs' reasoning capabilities to approximate and transfer existing knowledge rather than learning everything from scratch
\end{itemize}

\subsection{Research Questions}

This dissertation addresses three interconnected research questions toward realizing cost-effective wireless foundation models:

\textbf{RQ1: Can vision-based, self-supervised learning effectively transfer world knowledge to the wireless domain?}
\begin{itemize}
    \item How can wireless signals be represented as images (e.g., spectrograms, constellation diagrams) to leverage pre-trained vision models like ViT and YOLO?
    \item What self-supervised objectives (e.g., masked autoencoding) are most effective for learning rich, transferable representations from these signal-images?
    \item Can multimodal learning, combining different image representations (e.g., IQ samples and constellation diagrams), improve feature quality and robustness?
    \item How effectively do features learned by vision transformers transfer to wireless-specific tasks like modulation classification and signal detection?
    \item To what extent can this cross-modal approach reduce the dependency on large labeled datasets in the wireless domain?
\end{itemize}

\textbf{RQ2: Can LLMs enable open-set classification of wireless signals through world knowledge transfer?}
\begin{itemize}
    \item Can statistical signal features (cumulants, spectral properties) be converted to natural language prompts for LLM reasoning?
    \item Do pre-trained LLMs' in-context learning capabilities transfer to wireless signal classification without fine-tuning?
    \item What bridges the semantic gap between quantitative signal properties and linguistic descriptions?
    \item How competitive is training-free LLM-based classification compared to extensively trained task-specific models?
    \item Can LLMs provide interpretable reasoning for wireless classification decisions, enabling human-understandable explanations?
\end{itemize}

\textbf{RQ3: Can meta-learning enable rapid adaptation to new tasks with minimal labeled data?}
\begin{itemize}
    \item What episodic training strategies maximize generalization for wireless classification tasks?
    \item Can task-aware regularization within meta-learning prevent overfitting while maintaining adaptability?
    \item How can we design episode generation strategies that account for signal-specific characteristics like SNR variation and modulation similarity?
    \item What design principles enable efficient knowledge transfer across different channel conditions and modulation schemes?
    \item How many labeled examples are minimally required for effective adaptation to new wireless tasks?
\end{itemize}

\section{Scope}
\label{sec:scope}

This work focuses on automatic modulation classification (AMC) as the primary testbed for developing foundational principles for wireless physical-layer foundation models. The scope is defined by the following domain parameters and limitations.

\subsection{Primary Domain: Wireless Communication}

Our core application domain is automatic modulation classification under challenging conditions:
\begin{itemize}
    \item \textbf{Signal Representations:} We explore time-domain IQ samples, constellation diagrams, and noise patterns as complementary modalities
    \item \textbf{Operating Conditions:} SNR ranges from -10 dB to 30 dB, covering severe noise to clean signal regimes\cite{}
    \item \textbf{Modulation Schemes:} Focus on common digital modulations (BPSK, QPSK, 8PSK, QAM variants, FSK, GFSK)\cite{} with extensions to emerging schemes
    \item \textbf{Channel Models:} Primarily AWGN and fading channels\cite{} representative of practical wireless systems
    \item \textbf{Datasets:} %RadioML2016.10a/b\cite{}, RadioML2018.01a\cite{}, and 
    synthetic datasets for controlled experiments
\end{itemize}

\subsection{Scope Limitations}

While this work establishes foundations for wireless AI, certain aspects remain future work:
\begin{itemize}
    \item \textbf{End-to-End Communication:} Focus is on physical layer modulation classification, not complete communication system design\cite{}
    \item \textbf{Hardware Implementation:} Validation uses simulated and standard datasets rather than custom hardware deployments\cite{}
    \item \textbf{Adversarial Robustness:} While noise robustness is central, adversarial attacks\cite{} and countermeasures are not extensively explored
    \item \textbf{Multi-User/MIMO Scenarios:} Primarily single-user, single-antenna scenarios; multi-antenna\cite{} and interference management are secondary
    \item \textbf{Foundation Model Scale:} Due to computational constraints, models are smaller than vision/language foundation models\cite{}; principles established here enable future scaling
\end{itemize}

\section{Contributions} \label{sec:contributions}

This dissertation makes four primary technical contributions that collectively establish a modular, cost-effective pathway toward Wireless Physical-Layer Foundation Models (WPFMs). The first three contributions form the core components of the proposed framework, while the fourth provides a crucial benchmark and methodology for cross-modal knowledge transfer.

\subsection{Toward Wireless Foundation Models: A Modular, Cost-Effective Framework}
\label{subsec:unified_framework}

\textbf{Central Thesis Contribution:} We propose a modular and cost-effective framework for Wireless Physical-Layer Foundation Models (WPFMs) that strategically approximates and integrates world knowledge, aiming to demonstrate that:

\begin{itemize}
    \item \textbf{Self-Supervised Learning Component:} Multimodal masked autoencoders (DenoMAE) will learn universal wireless signal representations through unsupervised learning on unlabeled data, dramatically reducing labeled data requirements (80\% reduction) while approximating world knowledge of signal patterns
    \item \textbf{LLM Integration Component:} Large language models can perform open-set wireless classification by reasoning over compact, retrieval-augmented symbolic prompts rather than raw signal data. Plug-and-Play first established training-free feasibility; DiSC-AMC~\cite{disc_amc} then resolved its token-efficiency bottleneck, letting a fine-tuned 7B model match or surpass dedicated supervised encoders (83.0\% in-distribution, 82.5\% out-of-distribution) and a training-free 7B model outperform a 200B-parameter baseline at a fraction of the prompt length (71.9\% at 0.5K tokens vs.\ 69.92\% at 2.9K)
    \item \textbf{Hybrid Architecture Component:} DenoMAE 's hybrid generative-discriminative learning combined with an improved plug-and-play discriminative classifier will ensure robust performance across SNR conditions (-10 dB to 30+ dB), addressing the critical SNR imbalance problem through rich representation learning and condition-aware classification
    \item \textbf{Meta-Learning Component:} Meta-Task learning will enable rapid adaptation to novel modulation types with minimal labeled data (5-10 examples), efficiently applying integrated knowledge to new tasks and channel conditions
    \item \textbf{Cost-Effectiveness Principle:} The key to economical wireless foundation models lies in this strategic, modular framework design that leverages existing world knowledge rather than massive isolated training
\end{itemize}

This modular approach will provide both a conceptual blueprint and practical pathway toward next-generation wireless AI systems that are scalable, interpretable, and economically deployable. The framework's modularity will enable: (1) independent component optimization and updates; (2) targeted improvements without full system retraining; (3) flexible integration of future advances in self-supervised learning, LLMs, or meta-learning; (4) cost-effective deployment through training-free operation and minimal data requirements.

\section{Structure}
\label{sec:structure}

This dissertation is organized to guide readers from foundational concepts through technical contributions toward the unifying vision of wireless foundation models. The thesis follows a standard academic structure with nine chapters, each building progressively on previous content.

\begin{itemize}
    % \item \textbf{Chapter~\ref{chap:intro}: Introduction:} Establishes the vision of wireless foundation models as the organizing principle for this work. Articulates challenges in wireless communication that motivate foundation model approaches. Outlines five interconnected research questions (RQ1-RQ5), presents three main technical contributions (DenoMAE, Plug-and-Play, Meta-Task), and provides roadmap emphasizing the novel integration of self-supervised learning, LLM-based reasoning, and meta-learning for wireless signals.

    % \item \textbf{Chapter~\ref{chap:background}: Background and Related Work:} Provides comprehensive background consolidating literature across all technical areas. Reviews foundation models in computer vision and NLP, automatic modulation classification approaches, transformer architectures for wireless signals, self-supervised learning paradigms (contrastive learning, masked autoencoders), LLM capabilities and in-context learning, meta-learning frameworks (MAML, Prototypical Networks), and cross-modal learning. Establishes theoretical foundations and positions our contributions within the broader research landscape. Documents the three datasets used throughout this work (DenoMAE 2.0, Custom 10K, Drone Detection).

    % \item \textbf{Chapter~\ref{chap:methodology}: Methodology:} Describes the overall research methodology and framework connecting the three contributions. Presents the modular framework design showing how DenoMAE, Plug-and-Play, and Meta-Task components work independently and can integrate synergistically. Provides detailed descriptions of the three datasets used across all experiments, emphasizing experimental design, data preprocessing pipelines, evaluation metrics, and validation strategies. Establishes the methodological foundation for the empirical work in subsequent chapters.

    \item 	\textbf{Chapter~\ref{chap:denomae} and~\ref{chap:cdrf}: Self-Supervised Representation Learning (DenoMAE \& CDRF):} Presents Contribution 1, addressing RQ1, RQ4, and RQ5. Introduces DenoMAE, a multimodal masked autoencoder that treats noise as an informative modality. Details the architecture processing time-domain IQ samples and constellation diagrams through vision transformers. This chapter also introduces the CDRF benchmark, demonstrating the successful application of computer vision models to RF drone detection by representing signals as images. This combined work establishes that vision transformers can approximate world knowledge for wireless signals, enabling significant data reduction and cross-modal understanding.

    \item \textbf{Chapter~\ref{chap:plugandplay}: LLM-Based Reasoning for Open-Set AMC (Plug-and-Play and DiSC-AMC):} Presents Contribution 2, addressing RQ2. Introduces Plug-and-Play, which first demonstrates training-free wireless classification through LLM reasoning over statistical feature-to-language prompts, and diagnoses its token-efficiency bottleneck. Then introduces DiSC-AMC~\cite{disc_amc}, which resolves that bottleneck through feature discretization, self-supervised context retrieval, and candidate-space pruning, achieving competitive or superior accuracy relative to dedicated supervised encoders at a fraction of the token and parameter budget, both fine-tuned and training-free.

    \item 	\textbf{Chapter~\ref{chap:meta_task}: Meta-Learning for Rapid Adaptation (Meta-Task):} Presents Contribution 3, addressing RQ3. Introduces Meta-Task, which reformulates regularization as a learnable task within few-shot learning. Details the Task-Decoder architecture and demonstrates significant performance gains and faster convergence, establishing that meta-learning enables rapid adaptation to new wireless tasks with minimal labeled data.
    
    % \item 	\textbf{Chapter~\ref{chap:results}: Experimental Results and Analysis:} Consolidates and compares results across all three contributions. Presents comprehensive experimental validation including ablation studies, cross-dataset evaluation, and comparative analysis. Analyzes how each method performs across different SNR conditions, dataset sizes, and task complexities. Discusses performance trade-offs and computational efficiency. Evaluates the potential for component integration and synergies within the modular framework.
    
    \item 	\textbf{Chapter~\ref{chap:concl}: Conclusions and Future Work:} Synthesizes key findings and answers the three research questions, including a bounded, partly-affirmative answer to RQ3 grounded in a null result from testing Meta-Task directly on wireless AMC. States explicit authorship and role for each collaborative work used as evidence (Section~\ref{sec:authorship}), distinguishing first-author contributions (Plug-and-Play, DiSC-AMC, Meta-Task, shared-first-author CDRF) from co-authored representations used as evidence (DenoMAE, DenoMAE2.0). Discusses remaining limitations and outlines concrete future research directions, including a first (negative) fusion result, vision-language reasoning, and scoped open gaps identified during evaluation.
    
\end{itemize}